\subsection{Translation and protein evidence for predicted ORFs in GM12878}
% Draft (2026-10-21). Numbers: genome_annotation/results/gm12878_riboseq_*.tsv, gm12878_ms_*.tsv;
% notes/06_gm12878_three_assemblies.md.

We tested the GM12878 predictions against two independent experimental data sets from the same cell
line: ribosome profiling \citep{Cenik2015} and a deep, six-protease proteome \citep{Sinitcyn2023}.
Each distinct predicted CDS was classed by its relation to GENCODE: exact CDS match, compatible
(a correct partial CDS), other overlap with an annotated CDS, or unannotated (no overlap with any
annotated CDS).

\paragraph{Ribosome footprints.}
Predictions matching annotated CDSs were broadly covered by footprints. Of exact-match predictions,
83--89\% had at least ten P-sites, and 53--59\% passed a test for frame-0 enrichment. That test is
conservative here, because the RNase~A/T1 digestion used for these libraries gives weak 3-nt
periodicity (frame-0 fraction about 0.55 for the read lengths used). Unannotated predictions were
almost never covered: 4.3\%, 1.0\% and 0.6\% (LRAA, StringTie, Trinity), with 0.2--1.5\% passing
the frame test. The difference held when we compared transcripts with similar RNA abundance. At
5--20 TPM, 89\% of exact-match predictions were covered against 3.4\% of unannotated ones, so low
expression does not explain the lack of footprints.

\paragraph{Mass-spectrometry peptides.}
Of 758,603 peptides identified at 1\% FDR, 92--94\% of exact-match predictions of at least 100~aa
carried at least one peptide, nearly all of them two or more. Short exact-match predictions
(\textless100~aa) had 81--87\% coverage. Only 10--25\% of unannotated predictions shared any
peptide with the identified set, and these were almost entirely paralogous or pseudogene copies of
known protein-coding genes (for example \textit{NCF1B/C}, \textit{HERC2P3} and processed
ribosomal-protein pseudogenes). There, peptides cannot distinguish the copy from its parent gene.
Peptides absent from both GENCODE and UniProt that matched unannotated predictions were nearly all
GM12878 variant-allele peptides of parent genes mapping to duplicated copies, for example a
\textit{PWP2} copy on 21p and the \textit{NAPSB} pseudogene. The one credible exception was
\textit{C3orf86P}, supported by three independent peptides from three proteases. It was annotated
as protein-coding in the Ensembl release used for the original search and is a transcribed unitary
pseudogene in GENCODE v50. Because the original search database did not include our predictions,
this analysis confirms known proteins but cannot discover new ones.

\paragraph{Interpretation.}
Ribosome profiling and proteomics agree. Predictions at annotated coding loci are translated and
detected as protein, while the short unannotated ORFs gained at a 50-aa minimum show no evidence of
translation or protein in this cell line. Their apparent protein support comes from paralogs and
pseudogenes. Such predictions are best treated as candidates for evidence-based confirmation, not
as proteins.
\todo{Results for the homology-restriction series (clade-excluded Swiss-Prot); link to the
minimum-length sweep; candidate neo-ORFs supported only by transposon homology.}
